\section{Motivating Examples}\label{sec:examples}

The enum example in Sect.~\ref{sec:intro} is the headline case:
mutation to a parent value invalidates a reference still in scope
below it. Two further examples make the design constraints concrete.

\paragraph{Order matters.}
Here are two fragments that differ only in the order of their last
two statements:
\begin{lstlisting}
let x = 0; let r2 = &mut x; let r1 = &mut x;       // (a)
*r2 = 5; *r1
let x = 0; let r2 = &mut x; let r1 = &mut x;       // (b)
*r1 = 5; *r2
\end{lstlisting}
Both fragments take two mutable references to \texttt{x}. Fragment
(a) writes through one of them and reads through the other; fragment
(b) does the same with the roles swapped. A static borrow check sees
the same set of borrows in scope at the same program points and has
to reject both. The runtime sees the actual sequence of operations
and can be more precise: the read after the write is safe in either
fragment, because the read does not modify state and the earlier
write through the other handle has not been challenged by a later
write. A small variant---\texttt{let r1 = r2;}---makes both handles
copies of one reference, and the question becomes uninteresting.

\paragraph{Two mutable arguments must not overlap.}
A function takes two mutable references. The caller must make sure
they refer to disjoint pieces of state, since otherwise the callee
cannot reason about either argument locally. Passing
\texttt{\&mut x.f} and \texttt{\&mut x.g} is fine---they touch
different fields. Passing \texttt{\&mut x} and \texttt{\&mut x.g} is
not---the second is reachable through the first. The runtime check is
substructural: it has to ask, for any two reference arguments, whether
one is reachable from the other through field projections,
indexing, and so on, along the actual paths the program took.

\medskip

The three examples look different on the surface: variant mutation,
mutable aliasing, and an overlap at a call boundary. But in every
case, one operation invalidates the path through which another
reference would read or write state. That observation is what
motivates the model. The next section describes the data structure
that records those paths; Sect.~\ref{sec:semantics} describes the
single rule that decides which references the next destructive event
invalidates.
